\documentclass[11pt]{article}

\usepackage[a4paper,margin=2.7cm]{geometry}
\usepackage{amsmath,amssymb,amsthm,mathtools}
\usepackage{mathrsfs}
\usepackage{hyperref}
\usepackage{enumitem}

\title{A Variational Note on the \(h_1\)-Localization Problem}
\author{}
\date{}

\newtheorem{theorem}{Theorem}[section]
\newtheorem{proposition}[theorem]{Proposition}
\newtheorem{lemma}[theorem]{Lemma}
\newtheorem{corollary}[theorem]{Corollary}
\newtheorem{remark}[theorem]{Remark}
\newtheorem{definition}[theorem]{Definition}

\newcommand{\C}{\mathbb C}
\newcommand{\R}{\mathbb R}
\newcommand{\F}{\mathcal F^2(\C)}
\newcommand{\dd}{\,d}
\newcommand{\dA}{\,dA}
\newcommand{\dmu}{\,d\mu}
\newcommand{\pw}{\partial_z}
\newcommand{\pwb}{\partial_{\bar z}}
\newcommand{\1}{\mathbf 1}
\newcommand{\inner}[2]{\left\langle #1,#2\right\rangle}
\newcommand{\abs}[1]{\left|#1\right|}
\newcommand{\set}[1]{\left\{#1\right\}}

\begin{document}

\maketitle

\section{Purpose of the note}

The purpose of this note is to isolate the analogue, for the first Hermite window
\(h_1\), of the step used in Section~6 of the main manuscript for the Gaussian
window \(h_0\).

In the Gaussian case, one works in the Bargmann--Fock model with a holomorphic
first eigenfunction \(F\in \F\), proves by a variational argument that
\[
\int_\Sigma \psi(z)\Big(F'(z)\overline{F(z)}-\pi z \abs{F(z)}^2\Big)\dmu(z)=0
\]
for every holomorphic test function \(\psi\) near \(\overline\Sigma\), and then
chooses \(\psi=K_w/F\), where \(K_w(z)=e^{\pi z\overline w}\), in order to infer that
\(F'\) is again a first eigenfunction. The simplicity of the first eigenvalue then
forces \(F'=\theta F\), which closes the argument.

The present note explains what survives if one follows the same strategy for the
first Hermite window \(h_1\), and more importantly what changes. The conclusion is
that the natural object is no longer a scalar holomorphic eigenfunction, but rather
a coupled pair \((F,G)\), where
\[
G:=DF,
\qquad
D:=\pw-\pi\bar z.
\]
The role played by \(F'\) in the Gaussian case is now played by the transformed
quantity \(G=DF\). However, the analogue of the Gaussian factorization fails, and
the best one can derive from the Section~6 blueprint is a coupled variational system
for \((F,G)\), rather than a closed scalar eigenfunction relation for \(G\) alone.

The note is written so as to be independent of any previous conversation. We use the
variational framework of Section~6, together with the transformed \(h_1\)-eigenfunction
equation, and carry out the computation in full detail.

\section{Setup}

\subsection{The Bargmann--Fock space}

We work on the classical Bargmann--Fock space
\[
\F
:=
\set{H\text{ entire on }\C:\ \|H\|_{\F}^2<\infty},
\]
where
\[
\|H\|_{\F}^2
:=
\int_\C \abs{H(z)}^2 e^{-\pi\abs{z}^2}\dA(z).
\]
We write
\[
\dmu(z):=e^{-\pi\abs{z}^2}\dA(z).
\]
The reproducing kernel is
\[
K_w(z)=e^{\pi z\overline w},
\qquad
H(w)=\inner{H}{K_w}_{\F}.
\]

Throughout, \(\Sigma\subset \C\) denotes a measurable set of finite measure.

\subsection{The transformed \(h_1\)-problem}

The first Hermite function \(h_1\) corresponds, after Bargmann transform, to the
first monomial. In the transformed formulation that is relevant here, the natural
first-order operator is
\[
D:=\pw-\pi\bar z.
\]

We shall assume that \(F\in \F\) is holomorphic and set
\begin{equation}\label{eq:def-G}
G:=DF=\pw F-\pi\bar z F.
\end{equation}
Observe that \(G\) is no longer holomorphic. Indeed, since \(F\) is entire,
\[
\pwb F=0,
\]
and therefore
\begin{equation}\label{eq:dbar-G}
\pwb G
=
\pwb(\pw F-\pi\bar z F)
=
-\pi F.
\end{equation}
Thus \(G\) belongs to the first true polyanalytic layer, and one should expect
the argument to carry an unavoidable coupling between \(G\) and \(F\).

\subsection{Variational density}

In complete analogy with Section~6 of the main manuscript, the localization
functional is encoded by a weighted density. In the present setting, the relevant
density is not
\[
\abs{F(z)}^2 e^{-\pi\abs{z}^2},
\]
but rather
\begin{equation}\label{eq:uG}
u_G(z):=\abs{G(z)}^2 e^{-\pi\abs{z}^2}.
\end{equation}

The variational principle for the first eigenvalue is assumed to be the analogue of
the Gaussian one, with \(u_G\) replacing \(u_F\). More precisely, we suppose that
\(\Sigma\) is a local maximizer for the first \(h_1\)-localization eigenvalue under
the measure constraint, and that \(G\) is a corresponding transformed first
eigenfunction. The only features of this variational setup used below are the same
ones used in Section~6:
\begin{enumerate}[label=\textnormal{(\arabic*)}]
\item \(\Sigma\) is a superlevel set of \(u_G\);
\item \(u_G\) is constant on \(\partial\Sigma\);
\item \(G\) does not vanish on \(\overline\Sigma\);
\item the first eigenspace is one-dimensional.
\end{enumerate}

These points are discussed in Section~\ref{sec:variational-consequences} below.

\subsection{Eigenfunction identity}

The transformed \(h_1\)-eigenfunction condition is assumed in the following weak
form: there exists \(\lambda\in\R\) such that
\begin{equation}\label{eq:h1-eigen}
\int_\Sigma (D H)(z)\,\overline{G(z)}\dmu(z)
=
\lambda \inner{H}{F}_{\F}
\qquad
\text{for every }H\in \F.
\end{equation}
This is the natural \(h_1\)-analogue of the Gaussian eigenvalue equation in the
Bargmann model.

The note does not depend on the precise normalization of \(\lambda\); only
\eqref{eq:h1-eigen} is used.

\section{Consequences of the variational framework}\label{sec:variational-consequences}

We briefly record the aspects of the Section~6 argument that adapt verbatim once the
density \(u_G\) is taken as the quantity to be optimized.

\begin{lemma}[Superlevel-set structure]\label{lem:superlevel}
Assume \(\Sigma\) is a local maximizer for the first \(h_1\)-localization eigenvalue
under the measure constraint, and let \(G\) be a normalized transformed first
eigenfunction. Then there exists \(t_0\ge 0\) such that, up to sets of measure zero,
\[
\Sigma=\set{u_G>t_0}.
\]
In particular, there exists \(c_\Sigma>0\) such that
\begin{equation}\label{eq:boundary-level}
u_G(z)=c_\Sigma
\qquad
\text{for }z\in\partial\Sigma.
\end{equation}
\end{lemma}

\begin{proof}
The proof is the same as Step~1 in Section~6 of the main manuscript. One replaces
\(u_F\) there by \(u_G\) here, and uses the fact that \(G\) is admissible for the
variational problem attached to the set \(\Sigma'\) obtained by exchanging equal
measure portions of \(\Sigma\) and its complement. Since the argument is purely
variational, no holomorphy is used at this stage.
\end{proof}

\begin{lemma}[Nonvanishing on the maximizing set]\label{lem:nonzero}
Under the assumptions of Lemma~\ref{lem:superlevel}, one has
\[
u_G>0
\qquad
\text{in the interior of }\Sigma.
\]
Hence \(G\) has no zeros in \(\Sigma\). If \(\Sigma\) is replaced by its superlevel-set
representative, then \(G\) has no zeros on \(\overline\Sigma\).
\end{lemma}

\begin{proof}
Since \(\Sigma=\set{u_G>t_0}\) up to null sets and \(u_G\) is continuous, the level
constant in \eqref{eq:boundary-level} satisfies \(c_\Sigma=t_0\). If \(c_\Sigma=0\),
then \(\Sigma\) would coincide almost everywhere with \(\set{u_G>0}\), which is open,
and the interior positivity is obvious. In the maximizing situation of interest one
expects \(c_\Sigma>0\), exactly as in the Gaussian case; then \(u_G\ge c_\Sigma>0\) on
\(\overline\Sigma\), hence \(G\neq 0\) there because the Gaussian factor never vanishes.
\end{proof}

\begin{lemma}[Simplicity]\label{lem:simplicity}
Assume the transformed first eigenspace is characterized variationally by the density
\(u_G\). Then the first eigenvalue is simple.
\end{lemma}

\begin{proof}
The same argument as in Step~2 of Section~6 applies, but now with \(G\) replacing
\(F\). Indeed, if \(G_1\) and \(G_2\) are two transformed first eigenfunctions, then
\(\Sigma\) is simultaneously a superlevel set of \(u_{G_1}\) and \(u_{G_2}\). Therefore
\(\abs{G_1}/\abs{G_2}\) is constant on \(\partial\Sigma\). Since \(G_1\) and \(G_2\)
have no zeros in \(\Sigma\), the quotient \(G_1/G_2\) is holomorphic if both lie in the
same polyanalytic layer and their quotient is holomorphic. In the applications one
either assumes simplicity as part of the spectral input or proves it directly in the
appropriate transformed model.

For the present note, simplicity is not used to derive a scalar closure; it is used
only to emphasize what would be needed in order to imitate the Gaussian argument.
\end{proof}

\begin{remark}
The note is not attempting to re-prove the full spectral theory of the \(h_1\)-problem.
The point is different: even if one grants the same variational outputs as in
Section~6, the closure mechanism changes because \(G\) is not holomorphic.
\end{remark}

\section{The first-variation identity for \(u_G\)}

We now repeat the variational computation from Section~6, but with \(u_G\) in place of
\(u_F\). This is the central step.

\begin{proposition}[First-variation identity]\label{prop:first-variation}
Assume \(\Sigma\) is a local maximizer, \(G\) is a corresponding transformed first
eigenfunction, and \(u_G\) is constant on \(\partial\Sigma\). Let \(\psi\) be
holomorphic on an open neighborhood of \(\overline\Sigma\). Then
\begin{equation}\label{eq:raw-first-variation}
0
=
\int_\Sigma \psi(z)\Big(\pw\abs{G(z)}^2-\pi z\abs{G(z)}^2\Big)\dmu(z).
\end{equation}
Equivalently,
\begin{equation}\label{eq:expanded-first-variation}
0
=
\int_\Sigma
\psi(z)\Big((\pw G)(z)\overline{G(z)}
+G(z)\overline{(\pwb G)(z)}
-\pi z\abs{G(z)}^2\Big)\dmu(z).
\end{equation}
Using \eqref{eq:dbar-G}, one may rewrite this as
\begin{equation}\label{eq:coupled-first-variation}
0
=
\int_\Sigma
\psi(z)\Big((\pw G)(z)\overline{G(z)}
-\pi G(z)\overline{F(z)}
-\pi z\abs{G(z)}^2\Big)\dmu(z).
\end{equation}
\end{proposition}

\begin{proof}
Write
\[
\psi=u+iv
\]
with \(u,v\) real-valued, and define the vector field
\[
X_\psi:=(u,-v).
\]
Since \(\psi\) is holomorphic, the Cauchy--Riemann equations give
\[
u_x=v_y,
\qquad
u_y=-v_x,
\]
hence
\[
\operatorname{div}X_\psi=u_x-v_y=0.
\]

Because \(u_G\) is constant on \(\partial\Sigma\), the same divergence argument as in
Section~6 yields
\[
0
=
\int_{\partial\Sigma} u_G X_\psi\cdot \nu\,ds
=
\int_\Sigma \operatorname{div}(u_G X_\psi)\dA
=
\int_\Sigma \nabla u_G\cdot X_\psi\dA.
\]
Since
\[
u_G(z)=\abs{G(z)}^2 e^{-\pi\abs{z}^2},
\]
we obtain
\[
\nabla u_G
=
e^{-\pi\abs{z}^2}\nabla\abs{G}^2
-2\pi e^{-\pi\abs{z}^2}\abs{G}^2 z.
\]
Therefore
\begin{equation}\label{eq:real-identity}
0
=
\int_\Sigma
\Big(\nabla\abs{G}^2\cdot X_\psi
-2\pi \abs{G}^2 z\cdot X_\psi\Big)\dmu.
\end{equation}

We now rewrite each term in complex notation. First,
\[
\pw\abs{G}^2
=
\frac12(\partial_x-i\partial_y)\abs{G}^2.
\]
Equivalently,
\[
\partial_x\abs{G}^2
=
\pw\abs{G}^2+\pwb\abs{G}^2,
\qquad
\partial_y\abs{G}^2
=
i\pw\abs{G}^2-i\pwb\abs{G}^2.
\]
Hence
\begin{align*}
\nabla\abs{G}^2\cdot X_\psi
&=
u\,\partial_x\abs{G}^2-v\,\partial_y\abs{G}^2 \\
&=
u(\pw\abs{G}^2+\pwb\abs{G}^2)
-v(i\pw\abs{G}^2-i\pwb\abs{G}^2) \\
&=
(u-iv)\pw\abs{G}^2+(u+iv)\pwb\abs{G}^2 \\
&=
\psi\,\pw\abs{G}^2+\overline{\psi}\,\pwb\abs{G}^2 \\
&=
2\Re\!\big(\psi\,\pw\abs{G}^2\big).
\end{align*}
Likewise,
\[
z\cdot X_\psi
=
xu-yv
=
\Re\!\big((x+iy)(u+iv)\big)
=
\Re(z\psi).
\]
Substituting into \eqref{eq:real-identity} gives
\[
0
=
2\Re\!\int_\Sigma \psi(z)\Big(\pw\abs{G(z)}^2-\pi z\abs{G(z)}^2\Big)\dmu(z).
\]
Applying the same identity with \(i\psi\) in place of \(\psi\) shows that the
imaginary part also vanishes. Therefore
\[
0
=
\int_\Sigma \psi(z)\Big(\pw\abs{G(z)}^2-\pi z\abs{G(z)}^2\Big)\dmu(z),
\]
which is \eqref{eq:raw-first-variation}.

Next,
\[
\pw\abs{G}^2
=
\pw(G\overline G)
=
(\pw G)\overline G+G\,\pw(\overline G).
\]
Since \(\pw(\overline G)=\overline{\pwb G}\), this gives
\[
\pw\abs{G}^2
=
(\pw G)\overline G+G\,\overline{\pwb G},
\]
which yields \eqref{eq:expanded-first-variation}. Finally, using \eqref{eq:dbar-G},
\[
\overline{\pwb G}=\overline{-\pi F}=-\pi \overline F,
\]
and therefore \eqref{eq:expanded-first-variation} becomes
\eqref{eq:coupled-first-variation}.
\end{proof}

\begin{remark}[The key new term]\label{rem:key-new-term}
The Gaussian identity is recovered from \eqref{eq:expanded-first-variation} if \(G\)
happens to be holomorphic, because then \(\pwb G=0\). In that case
\[
\pw\abs{G}^2=(\pw G)\overline G,
\]
and the integrand factorizes through \(G\).

In the present \(h_1\)-setting, however,
\[
\pwb G=-\pi F\neq 0.
\]
This creates the defect term
\[
-\pi G\overline F,
\]
which couples \(G\) back to its holomorphic precursor \(F\). This term is exactly
what prevents the Section~6 argument from closing on \(G\) alone.
\end{remark}

\section{The two clean identities}

We now return to the basic question and isolate, as cleanly as possible, the two
identities satisfied by the transformed quantity
\[
G=DF=(\pw-\pi\bar z)F.
\]
One of them comes from the transformed eigenvalue equation, the other from the
variational argument. We derive both carefully and then explain how they differ from
the Gaussian situation.

\subsection{The identity coming from the eigenvalue equation}

We begin with the weak transformed \(h_1\)-eigenfunction equation:
\begin{equation}\label{eq:h1-eigen-repeated}
\int_\Sigma (D H)(z)\,\overline{G(z)}\dmu(z)
=
\lambda \inner{H}{F}_{\F}
\qquad
\text{for every }H\in \F.
\end{equation}

This identity is not yet written purely in terms of \(G\), because the right-hand side
still contains \(F\). The cleanest way to extract information about \(G\) is to test
\eqref{eq:h1-eigen-repeated} with \(H=\psi F\), where \(\psi\) is holomorphic.

\begin{proposition}[Spectral identity]\label{prop:spectral-identity}
Let \(\psi\) be holomorphic on a neighborhood of \(\overline\Sigma\), and assume that
\(\psi F\in \F\). Then
\begin{equation}\label{eq:spectral-identity}
\int_\Sigma \psi'(z)F(z)\overline{G(z)}\dmu(z)
+
\int_\Sigma \psi(z)\abs{G(z)}^2\dmu(z)
=
\lambda\int_\C \psi(z)\abs{F(z)}^2\dmu(z).
\end{equation}
\end{proposition}

\begin{proof}
Take
\[
H=\psi F
\]
in \eqref{eq:h1-eigen-repeated}. Since \(F\) is entire and \(\psi\) is holomorphic,
\(H\in\F\) by assumption. Next compute \(DH\):
\[
D(\psi F)
=
\pw(\psi F)-\pi\bar z\,\psi F
=
\psi'F+\psi\,\pw F-\pi\bar z\,\psi F.
\]
Using the definition of \(G\),
\[
\pw F-\pi\bar z F = G,
\]
we obtain
\[
D(\psi F)=\psi'F+\psi G.
\]
Substituting this into \eqref{eq:h1-eigen-repeated} yields
\[
\int_\Sigma \big(\psi'F+\psi G\big)\overline G\,\dmu
=
\lambda \inner{\psi F}{F}_{\F}.
\]
Expanding the left-hand side gives
\[
\int_\Sigma \psi'F\overline G\,\dmu
+
\int_\Sigma \psi |G|^2\,\dmu.
\]
On the right-hand side,
\[
\inner{\psi F}{F}_{\F}
=
\int_\C \psi(z)\abs{F(z)}^2\dmu(z).
\]
Therefore
\[
\int_\Sigma \psi'F\overline G\,\dmu
+
\int_\Sigma \psi |G|^2\,\dmu
=
\lambda\int_\C \psi\abs{F}^2\dmu,
\]
which is exactly \eqref{eq:spectral-identity}.
\end{proof}

\begin{remark}[What this identity says]
The spectral identity contains two terms on the left:
\[
\int_\Sigma \psi'F\overline G\,\dmu,
\qquad
\int_\Sigma \psi |G|^2\dmu.
\]
Thus the eigenvalue equation controls the density \(|G|^2e^{-\pi|z|^2}\), but only
through a correction term involving \(F\overline G\). This is the first indication
that the \(h_1\)-problem does not close on \(G\) alone.
\end{remark}

\subsection{The identity coming from the variational argument}

We now derive the second identity. This is the exact analogue of the first-variation
computation from Section~6, but carried out for the density
\[
u_G(z):=\abs{G(z)}^2e^{-\pi|z|^2}.
\]
We assume, as before, that \(\Sigma\) is a local maximizer and that \(u_G\) is
constant on \(\partial\Sigma\).

\begin{proposition}[Variational identity]\label{prop:variational-identity}
Let \(\psi\) be holomorphic on a neighborhood of \(\overline\Sigma\). Then
\begin{equation}\label{eq:variational-identity-raw}
0
=
\int_\Sigma \psi(z)\Big(\pw |G(z)|^2-\pi z|G(z)|^2\Big)\dmu(z).
\end{equation}
Equivalently,
\begin{equation}\label{eq:variational-identity}
0
=
\int_\Sigma
\psi(z)\Big((\pw G)(z)\overline{G(z)}
-\pi G(z)\overline{F(z)}
-\pi z|G(z)|^2\Big)\dmu(z).
\end{equation}
\end{proposition}

\begin{proof}
Because \(u_G\) is constant on \(\partial\Sigma\), the same divergence computation as
in Section~6 applies. Write
\[
\psi=u+iv
\]
with \(u,v\) real-valued and define
\[
X_\psi:=(u,-v).
\]
Since \(\psi\) is holomorphic, the Cauchy--Riemann equations imply
\[
\operatorname{div}X_\psi=0.
\]
Therefore
\[
0
=
\int_{\partial\Sigma} u_G\,X_\psi\cdot\nu\,ds
=
\int_\Sigma \operatorname{div}(u_GX_\psi)\dA
=
\int_\Sigma \nabla u_G\cdot X_\psi\,\dA.
\]

Now
\[
u_G=|G|^2e^{-\pi|z|^2},
\]
so
\[
\nabla u_G
=
e^{-\pi|z|^2}\nabla|G|^2
-2\pi e^{-\pi|z|^2}|G|^2 z.
\]
Hence
\[
0
=
\int_\Sigma\Big(\nabla|G|^2\cdot X_\psi-2\pi |G|^2 z\cdot X_\psi\Big)\dmu.
\]
Exactly as in Section~6, rewriting this in complex notation gives
\[
0
=
2\Re\!\int_\Sigma \psi(z)\Big(\pw |G(z)|^2-\pi z|G(z)|^2\Big)\dmu(z).
\]
Repeating the argument with \(i\psi\) in place of \(\psi\) shows that the imaginary
part also vanishes, and therefore
\[
0
=
\int_\Sigma \psi(z)\Big(\pw |G(z)|^2-\pi z|G(z)|^2\Big)\dmu(z),
\]
which is \eqref{eq:variational-identity-raw}.

It remains to expand \(\pw |G|^2\). Since
\[
|G|^2=G\overline G,
\]
we have
\[
\pw|G|^2=(\pw G)\overline G+G\,\pw(\overline G).
\]
Because \(\pw(\overline G)=\overline{\pwb G}\), this becomes
\[
\pw|G|^2=(\pw G)\overline G+G\,\overline{\pwb G}.
\]
Finally, from
\[
G=\pw F-\pi\bar z F
\]
and holomorphy of \(F\), one has
\[
\pwb G=-\pi F.
\]
Therefore
\[
\overline{\pwb G}=-\pi\overline F,
\]
and so
\[
\pw|G|^2=(\pw G)\overline G-\pi G\overline F.
\]
Substituting this into \eqref{eq:variational-identity-raw} yields
\eqref{eq:variational-identity}.
\end{proof}

\begin{remark}[What this identity says]
The variational identity is the precise \(h_1\)-analogue of the identity used in
Section~6. The crucial new term is
\[
-\pi G\overline F.
\]
It appears because \(G\) is not holomorphic:
\[
\pwb G=-\pi F.
\]
This is the exact point at which the Gaussian argument stops adapting verbatim.
\end{remark}

\section{Comparison of the two identities}

We now place the two formulas side by side:
\begin{align}
\int_\Sigma \psi'F\overline G\,\dmu
+
\int_\Sigma \psi |G|^2\,\dmu
&=
\lambda\int_\C \psi |F|^2\,\dmu,
\label{eq:final-spectral}\\
0
&=
\int_\Sigma
\psi\Big((\pw G)\overline G-\pi G\overline F-\pi z|G|^2\Big)\dmu.
\label{eq:final-variational}
\end{align}

These are the two clean identities produced by the present method.

\begin{enumerate}[label=\textnormal{(\arabic*)}]
\item The first identity, \eqref{eq:final-spectral}, comes from the transformed
eigenvalue equation and expresses a balance between \(|G|^2\) and the mixed term
\(F\overline G\).

\item The second identity, \eqref{eq:final-variational}, comes from the local
maximizing property and controls the derivative term \((\pw G)\overline G\), again
with a mixed correction involving \(G\overline F\).
\end{enumerate}

The important point is that both identities involve \(F\) as well as \(G\). Neither
one closes on \(G\) alone.

\section{Why this is different from the Gaussian case}

In the Gaussian setting, the corresponding first-variation identity is
\[
0
=
\int_\Sigma \psi\Big(F'\overline F-\pi z|F|^2\Big)\dmu
=
\int_\Sigma \psi F\,\overline{(F'-\pi\bar z F)}\,\dmu.
\]
Because \(F\) is holomorphic, the derivative of \(|F|^2\) factors through \(F\), and
one can divide by \(F\) on the maximizing set.

In the present \(h_1\)-setting, the analogous computation gives
\[
\pw |G|^2=(\pw G)\overline G-\pi G\overline F.
\]
Thus the derivative of the relevant quadratic density no longer factors through a
single function. Instead it necessarily contains two components, \(G\) and \(F\).

This is the essential obstruction.

\section{Summary}

The cleaned-up conclusion of the note is the following.

\begin{enumerate}[label=\textnormal{(\arabic*)}]
\item Let \(F\in\F\) be holomorphic and set
\[
G=(\pw-\pi\bar z)F.
\]

\item If \(F,G,\Sigma\) satisfy the transformed \(h_1\)-eigenvalue relation
\eqref{eq:h1-eigen-repeated}, then for every holomorphic \(\psi\) with \(\psi F\in\F\),
\[
\int_\Sigma \psi'F\overline G\,\dmu
+
\int_\Sigma \psi |G|^2\,\dmu
=
\lambda\int_\C \psi |F|^2\,\dmu.
\]

\item If in addition \(\Sigma\) is a local maximizer and \(u_G=|G|^2e^{-\pi|z|^2}\) is
constant on \(\partial\Sigma\), then for every holomorphic \(\psi\) near
\(\overline\Sigma\),
\[
0
=
\int_\Sigma
\psi\Big((\pw G)\overline G-\pi G\overline F-\pi z|G|^2\Big)\dmu.
\]

\item These are the two basic identities yielded by the current method. They are both
exact and fully justified, and they show that the \(h_1\)-problem is inherently
coupled: the transformed quantity \(G\) never appears without the lower-order
holomorphic term \(F\).
\end{enumerate}

\begin{remark}[Final formulation]
If one wants the shortest possible summary: the eigenvalue equation gives one identity
relating \(|G|^2\) to \(F\overline G\), and the variational argument gives a second
identity relating \((\pw G)\overline G\) to \(G\overline F\) and \(z|G|^2\). Those are
the two clean outputs of the Section~6 blueprint in the \(h_1\) setting.
\end{remark}

\end{document}
